\section{Introduction}
\label{sec:introduction_dynamic}

Adiabatic shear localization is usually introduced as a competition between work hardening, rate sensitivity and thermal softening. Classical perturbation theories sharpened this picture by adding inertia and thermal diffusion and by asking when a homogeneous simple-shear solution admits a growing disturbance \cite{Bai1982,BaiDodd1992,Wright2002,WrightBatra1985,MolinariClifton1987}. Subsequent analyses clarified thermomechanical coupling, loss of stability and the role of regularizing physics in pattern selection \cite{BenallalBigoni2004,ArriagaWaisman2017,JacqueyRattezVeveakis2021}. In their natural domain, these theories answer an indispensable question: which root of a scalar, frozen dispersion relation has the largest instantaneous real part? They do not, however, determine which perturbation of an anisotropic crystal accumulates the largest observable amplification over an interval during which the tangent, active slip systems, dislocation populations and temperature all evolve. This distinction is the solid-mechanics counterpart of nonmodal stability theory, in which the propagator---rather than one instantaneous eigenvalue---is the primary finite-time object \cite{Schmid2007}.

That distinction is physical rather than semantic. HCP crystals possess competing basal, prismatic and pyramidal slip systems whose resolved driving forces depend on orientation. Crystal-plasticity calculations show that slip competition, texture and microstructural heterogeneity strongly affect localization \cite{FranzAbedMeraimBerveiller2013,ZhangEakinsDunne2016,MouradEtAl2017,AnEtAl2023}. Thermomechanically coupled formulations add heat generation and evolving resistance \cite{LiRomeroSegurado2019,DaiSong2022}; micromorphic and strain-gradient formulations add internal fields, grain-boundary interactions, length-dependent recovery and FFT-based localization regularization \cite{Gurtin2002,Gurtin2008,VanBeersEtAl2013,AlipourReeseWulfinghoff2019,PhalkeEtAl2022,LameJouybariEtAl2024,LameJouybariEtAl2025}. The resulting perturbation operator is both high dimensional and nonautonomous. Its eigenvectors can rotate, its branches can exchange dominance, and non-normal interactions can amplify a disturbance that is not represented by the leading frozen eigenvector at any single time.

The unresolved question is therefore specific: \emph{for an orientation-dependent, multi-slip, dislocation--thermal HCP crystal whose base state continues to evolve, which initial disturbance obtains the largest cumulative gain over the complete loading history, and why do its direction, wavenumber and state composition change with the observation horizon?} Integrating the largest frozen growth rate is not an answer unless the evolving operators and their modal subspaces satisfy restrictive commutation and compatibility conditions. Those conditions are neither guaranteed by a positive frozen eigenvalue nor tested by a conventional dispersion plot.

There is also a structural obstacle. With displacement inertia, temperature and internal variables as differential coordinates but micromorphic slip fields as algebraic coordinates, the Fourier-linearized crystal problem is a quadratic differential--algebraic descriptor. Descriptor-system transient analysis requires the finite deflating subspace and its induced norm to be handled explicitly \cite{EmbreeKeeler2017}. Adding fictitious storage to the algebraic fields creates artificial fast modes; eliminating them without a determinant and reconstruction audit can silently delete or move physical roots. A defensible finite-time theory must first produce an ordinary generator that is exactly equivalent on the finite spectrum and only then define an observable propagator.

This paper resolves those two levels in sequence (Fig.~\ref{fig:ft_framework}). First, an 84-coordinate HCP descriptor is reduced exactly to a 66-coordinate quadratic pencil and a 69-state generator, with the eliminated 18-coordinate microfield reconstructed algebraically. Second, the chronological propagator is evaluated in a fixed physical scaling. A projected constitutive input--output gain defines the set-valued direction--wavenumber localization selector. The central frozen-spectrum test is deliberately separate: it compares the complete 69-state propagator norm with the full-state spectral abscissa on the same branch, base history and norm. This avoids assigning a full-state frozen rate to a projected subspace that is not dynamically invariant.

Six additional evidence layers are used. The generator is checked against an independent generalized Schur decomposition of the augmented descriptor. Stationary-normal, commuting-normal and invariant-projection controls test when frozen accumulation is legitimate. The discrimination is repeated for three independently selected crystal orientations, while commutators, modal condition numbers, eigenvector rotation and order-reversed propagation diagnose its origin. Eleven norm definitions and seven selector pairs test conditioning of the localization result. The direction--wavenumber basin is challenged by nested Sobol reconnaissance, explicit prior-branch anchors, independent refinements and 129-state verification; because the blind design misses the terminal basin, this is retained only as an anchor-assisted basin audit. Finally, the actual near-onset and terminal singular vectors are advanced in a declared four-mode nonlinear Galerkin replay. A separate dealiased 16/32/64-cell audit independently varies the spatial cutoff, time step and integration format and localizes the first raw all-mode failure; the retained state and gain converge, but the unresolved right-hand-side energy rejects complete nonlinear spectral closure.

The scope is deliberately bounded. The scalar Bai implementation is an independently tested classical anchor, but a strict analytical HCP-to-Bai limit is not claimed. The search audit demonstrates retention of known basins, not independent global discovery. The critical time is conditioned on the chosen initial time, norm, selectors and threshold, and the titanium-like parameter card is a literature-constrained verification seed rather than a specimen-identified material. The nonthermal power remainder is a passive ledger rather than an identified dislocation free energy. These boundaries allow the paper to make one strong claim: a crystal-resolved, nonautonomous propagator answers a perturbation-selection question that frozen theories cannot answer, while keeping nonlinear length-scale, microstructural-energy and propagation-resistance closure as open problems.

\begin{figure}[t]
\centering
\includegraphics[width=0.99\textwidth]{fig01_theory_variable_framework.pdf}
\caption{Theory and variable framework. The classical scalar Bai problem provides the instantaneous thermoplastic anchor; the HCP formulation resolves orientation, 18 slip systems, dislocation populations, temperature, inertia and algebraic micromorphic fields; exact descriptor reduction then permits time-ordered finite-time selection. The strict analytical HCP-to-Bai limit remains an explicit open gate.}
\label{fig:ft_framework}
\end{figure}
